\subsection{Interrupting Execution Flow Subsystems}

Loops normally follow a stable repetition pattern. A \texttt{while} loop returns to its condition. A \texttt{for} loop asks its iterator for the next value. This ordinary loop route can be interrupted from inside the loop body.

Python provides two local loop-control statements for this purpose:

\begin{itemize}
    \item \texttt{break}: leave the loop immediately;
    \item \texttt{continue}: skip the rest of the current iteration and move to the next loop decision.
\end{itemize}

These statements are not expressions. They do not produce values. They alter the control-flow route.

\begin{quote}
\texttt{break} exits the loop. \texttt{continue} exits only the current iteration.
\end{quote}

\subsubsection{Terminating the Iteration Invariant: The Immediate Exit Properties of \texttt{break}}

The \texttt{break} statement terminates the nearest enclosing loop immediately. When Python reaches \texttt{break}, it does not finish the remaining statements in the current loop body. It jumps to the first statement after the loop.

\begin{verbatim}
items = ["soup", "pretzels", "muffin", "coffee"]

for item in items:
    print(f"checking: {item}")

    if item == "muffin":
        print("muffin found, stopping loop")
        break

    print("not the muffin yet")

print("after loop")
\end{verbatim}

Possible output:

\begin{verbatim}
checking: soup
not the muffin yet
checking: pretzels
not the muffin yet
checking: muffin
muffin found, stopping loop
after loop
\end{verbatim}

The line:

\begin{verbatim}
print("not the muffin yet")
\end{verbatim}

does not execute for \texttt{"muffin"}, because \texttt{break} exits the loop before Python reaches that statement.

The route is:

\begin{verbatim}
get next item
    -> print checking
    -> test item == "muffin"
        true  -> print found
              -> break
              -> jump after loop
        false -> print not found
              -> ask for next item
\end{verbatim}

The \texttt{break} statement exits only the nearest loop:

\begin{verbatim}
groups = [
    ["Jerry", "Elaine"],
    ["Bart", "Lisa"],
]

for group in groups:
    print(f"group: {group}")

    for name in group:
        print(f"checking name: {name}")

        if name == "Bart":
            print("found Bart, leaving inner loop")
            break

    print("after inner loop")

print("after outer loop")
\end{verbatim}

Possible output:

\begin{verbatim}
group: ['Jerry', 'Elaine']
checking name: Jerry
checking name: Elaine
after inner loop
group: ['Bart', 'Lisa']
checking name: Bart
found Bart, leaving inner loop
after inner loop
after outer loop
\end{verbatim}

The \texttt{break} exits the inner \texttt{for name in group:} loop. It does not exit the outer \texttt{for group in groups:} loop.

The same rule applies to \texttt{while}:

\begin{verbatim}
n = 1

while True:
    print(f"n: {n}")

    if n == 3:
        print("stop at three")
        break

    n = n + 1

print("loop stopped")
\end{verbatim}

Possible output:

\begin{verbatim}
n: 1
n: 2
n: 3
stop at three
loop stopped
\end{verbatim}

Here \texttt{while True:} creates an intentionally open loop. The actual exit route is inside the body. Without the \texttt{break}, the loop would not stop by itself.

This pattern should be used carefully. A \texttt{while True:} loop is safe only if the programmer can clearly identify the internal exit path.

\begin{quote}
A loop containing \texttt{break} has an additional exit edge from inside the body to the statement after the loop.
\end{quote}

The \texttt{break} statement is often used when the loop is searching for something:

\begin{verbatim}
names = ["Jerry", "Elaine", "Kramer", "Newman"]
target = "Kramer"
found = False

for name in names:
    print(f"checking: {name}")

    if name == target:
        found = True
        print(f"found: {name}")
        break

print(f"found status: {found}")
\end{verbatim}

Possible output:

\begin{verbatim}
checking: Jerry
checking: Elaine
checking: Kramer
found: Kramer
found status: True
\end{verbatim}

The list object is the reusable traversal source:

\begin{verbatim}
names = ["Jerry", "Elaine", "Kramer", "Newman"]

print(f"type(names): {type(names)}")
print(f"dir(names) sample: {dir(names)[:10]}")
print(f"len(names): {len(names)}")
print(f"has __iter__: {'__iter__' in dir(names)}")
print(f"has __getitem__: {'__getitem__' in dir(names)}")
\end{verbatim}

Possible output:

\begin{verbatim}
type(names): <class 'list'>
dir(names) sample: ['__add__', '__class__', '__class_getitem__',
'__contains__', '__delattr__', '__delitem__', '__dir__', '__doc__',
'__eq__', '__format__']
len(names): 4
has __iter__: True
has __getitem__: True
\end{verbatim}

The \texttt{\_\_iter\_\_()} method lets the list produce an iterator for the loop. The \texttt{\_\_getitem\_\_()} method supports indexed access. The \texttt{break} statement does not modify the list. It modifies only the route of execution through the loop.

A loop stopped by \texttt{break} may leave some values unvisited:

\begin{verbatim}
nums = range(1, 10)

for n in nums:
    print(f"n: {n}")

    if n == 4:
        print("breaking at 4")
        break

print("done")
\end{verbatim}

Possible output:

\begin{verbatim}
n: 1
n: 2
n: 3
n: 4
breaking at 4
done
\end{verbatim}

The values \texttt{5}, \texttt{6}, \texttt{7}, \texttt{8}, and \texttt{9} are not requested from the traversal route. The loop ended early.

\subsubsection{Short-Circuiting Current Iterations: The Jump Mechanics of \texttt{continue}}

The \texttt{continue} statement does not exit the whole loop. It exits the current iteration only. When Python reaches \texttt{continue}, it skips the rest of the current loop body and returns to the loop's next decision point.

In a \texttt{for} loop, that means Python asks the iterator for the next value. In a \texttt{while} loop, that means Python re-evaluates the condition.

\begin{verbatim}
items = ["soup", "", "pretzels", "", "muffin"]

for item in items:
    if item == "":
        print("empty item skipped")
        continue

    print(f"processing: {item}")

print("loop finished")
\end{verbatim}

Possible output:

\begin{verbatim}
processing: soup
empty item skipped
processing: pretzels
empty item skipped
processing: muffin
loop finished
\end{verbatim}

When \texttt{item == ""}, the loop prints the skip message and reaches \texttt{continue}. The later line:

\begin{verbatim}
print(f"processing: {item}")
\end{verbatim}

is skipped for that iteration only. The loop then moves on to the next item.

The route is:

\begin{verbatim}
get next item
    -> test item == ""
        true  -> print skipped
              -> continue
              -> ask for next item
        false -> print processing
              -> ask for next item
\end{verbatim}

The \texttt{continue} statement is useful for guard-style filtering. It lets the loop reject cases early and keep the main processing path less nested.

Without \texttt{continue}:

\begin{verbatim}
items = ["soup", "", "pretzels"]

for item in items:
    if item != "":
        print(f"processing: {item}")
    else:
        print("empty item skipped")
\end{verbatim}

With \texttt{continue}:

\begin{verbatim}
items = ["soup", "", "pretzels"]

for item in items:
    if item == "":
        print("empty item skipped")
        continue

    print(f"processing: {item}")
\end{verbatim}

Both versions can produce the same result. The second version says: reject the unwanted case immediately, then let the rest of the body describe the normal case.

In a \texttt{while} loop, \texttt{continue} jumps back to the condition test. This means the progress update must still happen before the \texttt{continue}, or the loop may become infinite.

Safe version:

\begin{verbatim}
n = 0

while n < 5:
    n = n + 1

    if n == 3:
        print("skipping three")
        continue

    print(f"n: {n}")

print("done")
\end{verbatim}

Possible output:

\begin{verbatim}
n: 1
n: 2
skipping three
n: 4
n: 5
done
\end{verbatim}

The update:

\begin{verbatim}
n = n + 1
\end{verbatim}

happens before the possible \texttt{continue}. Therefore the loop still moves toward termination.

Dangerous version:

\begin{verbatim}
n = 0

while n < 5:
    if n == 0:
        print("stuck before progress")
        continue

    n = n + 1
\end{verbatim}

Here \texttt{n} starts as \texttt{0}. The condition \texttt{n == 0} is true. Python prints the message and reaches \texttt{continue}. The update is skipped. Then Python checks \texttt{n < 5} again, finds it still true, and repeats forever.

So the practical rule for \texttt{while} loops is:

\begin{quote}
If a \texttt{continue} can run before the loop's progress update, check carefully that the loop still has a route toward termination.
\end{quote}

The difference between \texttt{break} and \texttt{continue} is therefore exact:

\begin{verbatim}
break:
    leave the nearest loop completely

continue:
    leave the current iteration only
    continue with the next loop decision
\end{verbatim}

A final example shows both statements in the same loop:

\begin{verbatim}
items = ["soup", "", "pretzels", "stop", "muffin"]

for item in items:
    print(f"seen: {item!r}")

    if item == "":
        print("skip empty item")
        continue

    if item == "stop":
        print("stop marker found")
        break

    print(f"processed: {item}")

print("after loop")
\end{verbatim}

Possible output:

\begin{verbatim}
seen: 'soup'
processed: soup
seen: ''
skip empty item
seen: 'pretzels'
processed: pretzels
seen: 'stop'
stop marker found
after loop
\end{verbatim}

The empty string causes \texttt{continue}, so only that iteration is skipped. The string \texttt{"stop"} causes \texttt{break}, so the entire loop terminates. The later value \texttt{"muffin"} is never reached.

\begin{quote}
\texttt{continue} preserves the loop but shortens the current iteration. \texttt{break} destroys the current loop route entirely and resumes after the loop.
\end{quote}